\section{Glosario}

A continuación se encuentran los conceptos y acrónimos relevantes para el entendimiento del Kit Pagos Colombia:

\begin{itemize}
  \item \textbf{Pasarela de pago:} Una pasarela de pago es un sistema de software que actúa como intermediario entre un comercio y las entidades financieras, autorizando y verificando transacciones de pago en línea mediante la transmisión segura de datos entre el comprador, el comercio y el banco procesador. (Hisyam \& Manuaba, 2022)
  \item \textbf{API - Interfaz de programación de aplicaciones:} Una API es un contrato de comunicación entre sistemas de software que define cómo un programa puede solicitar servicios o datos a otro, sin necesidad de conocer su funcionamiento interno. En el contexto de pagos digitales, las APIs permiten que una aplicación comercial interactúe con una pasarela de pago para crear transacciones, consultar estados y recibir resultados de forma estandarizada.(Fielding, 2000)
  \item \textbf{Framework:} En ingeniería de software, un framework se entiende como una estructura compuesta por componentes reutilizables que permiten estandarizar y agilizar el desarrollo de aplicaciones, apoyándose en principios de reutilización y en el uso de estándares definidos (Amazon Web Services, s. f.).
  \item \textbf{Patrones de diseño:} Los patrones de diseño son soluciones recurrentes y documentadas a problemas comunes en el diseño de software. No son código reutilizable directamente, sino descripciones o plantillas que guían cómo estructurar componentes para resolver un problema específico en un contexto determinado. (\textit{Patrones de Diseño GoF Por Pagina | PDF}, s. f.).
  \item \textbf{GoF:} En ingeniería de software, GoF es un acrónimo de “Gang of Four”(Pandilla de cuatro) que hace referencia a los cuatro autores (Erich Gamma, Richard Helm, Ralph Johnson y John Vlissides) del libro de 1994, \textit{Design Patterns: Elements of Reusable Object-Oriented Software.} Definen 23 patrones de diseño óptimos para resolver problemas reiterados en la programación orientada a objetos. (\textit{Patrones de Diseño GoF Por Pagina | PDF}, s. f.).
  \item \textbf{Arquitectura hexagonal:} La arquitectura hexagonal, propuesta por Alistair Cockburn, es un estilo que organiza una aplicación alrededor de un núcleo de lógica de negocio aislado. Este núcleo se comunica mediante puertos y adaptadores que lo conectan con tecnologías externas como bases de datos, APIs o interfaces de usuario. Su objetivo es mantener el núcleo independiente de cualquier tecnología o proveedor externo. (\textit{Hexagonal architecture pattern - AWS Prescriptive Guidance}, s. f.)
  \item \textbf{Normalización de APIs:} La normalización de APIs es el proceso de unificar las respuestas, estructuras de datos y comportamientos de múltiples servicios externos bajo un formato común y consistente. En sistemas que integran proveedores heterogéneos, la normalización reduce el acoplamiento entre el consumidor y cada proveedor individual, facilitando la intercambiabilidad de servicios. (Taherkordi et al., 2011)
  \item \textbf{Entornos de prueba:} Un entorno sandbox en el contexto de APIs de pago es un ambiente de simulación aislado provisto por el proveedor que replica el comportamiento de su sistema productivo sin procesar transacciones reales. Permite a los desarrolladores verificar la correcta integración de sus sistemas durante la fase de desarrollo, reduciendo el riesgo de errores en producción. (Ammann \& Offutt, 2016)
  \item \textbf{Integración de sistemas:} La integración de sistemas es el proceso de conectar subsistemas o componentes independientes para que funcionen como un conjunto coordinado. En arquitecturas orientadas a servicios, la integración implica gestionar diferencias en protocolos, formatos de datos y modelos de comunicación entre sistemas que no fueron diseñados para interoperar directamente.(\textit{Hohpe, G. and Woolf, B. (2003) Enterprise Integration Patterns Designing, Building, and Deploying Messaging Solutions. Addison-Wesley. - References - Scientific Research Publishing}, s. f.)
  \item \textbf{Webhook}: Un webhook es un mecanismo de notificación asíncrona entre sistemas mediante el cual un proveedor envía una solicitud HTTP POST hacia una URL registrada por el receptor cuando ocurre un evento determinado. En el contexto de pasarelas de pago, los webhooks notifican el resultado final de una transacción una vez que el procesamiento bancario ha concluido, de forma independiente a la sesión HTTP original. (Hohpe \& Woolf, 2003).
  \item \textbf{SDK (Software Development Kit):} Un SDK es un conjunto de herramientas, librerías, documentación y ejemplos de código que permiten a los desarrolladores construir aplicaciones que interactúen con una plataforma o servicio específico. En el contexto de este proyecto, el SDK del Kit Pagos Colombia es la biblioteca principal que abstrae la lógica de integración con las pasarelas de pago. (Amazon Web Services, s. f.).
  \item \textbf{Métricas CK (Chidamber y Kemerer):} Conjunto de métricas orientadas a objetos propuesto por Chidamber y Kemerer (1994) para medir la complejidad, el acoplamiento y otras propiedades estructurales del código. Las tres métricas usadas en este proyecto son WMC (Weighted Methods per Class), CBO (Coupling Between Object Classes) y RFC (Response For a Class). (Chidamber \& Kemerer, 1994)
  \item \textbf{PSE (Pagos Seguros en Línea):} Botón de pago de ACH Colombia con el que el comprador paga por internet debitando su cuenta de ahorros o corriente, sin tarjeta. El comercio crea el pago en su pasarela, redirige al comprador al portal de su banco y conoce el resultado al final, por consulta o por webhook. (ACH Colombia, s. f.)
  \item \textbf{Tokenización:} Reemplazo del número de la tarjeta (PAN) por un valor sustituto, el token, que no permite recuperar el número original. El comercio guarda y transmite el token en lugar del PAN, lo que reduce los sistemas que entran al alcance de PCI DSS. En este proyecto, el módulo de navegador obtiene el token directamente de la pasarela, de modo que el número no pasa por el servidor del comercio. (PCI Security Standards Council, 2011)
  \item \textbf{DSR (Design Science Research):} Metodología de investigación utilizada en sistemas de información e ingeniería de software cuando el resultado central de la investigación es la construcción y evaluación de un artefacto que resuelve un problema real. Organiza el trabajo en seis fases: identificación del problema, definición de objetivos, diseño, desarrollo, demostración y comunicación. (Peffers et al., 2007).
\end{itemize}
